% Экспериментальная установка (main.tex)
\section{Экспериментальная установка}\label{sec:setup}

\paragraph{Модели.} Qwen3-0.6B и 1.7B в варианте Base. Обе модели используют классическое softmax-внимание с GQA,
QK-нормализацией и RoPE. Модели Qwen3.5 малых размеров не подходят:
в них три из четырёх слоёв "--- линейное внимание Gated DeltaNet, и
первые слои не содержат softmax-внимания. Мы не обучаем модели
заново, а берём опубликованные веса. Исходная модель (столбец <<Исх.>>)
проверяется по тем же бенчмаркам. Второе семейство "--- архитектура Llama
без QK-нормализации и со связанными эмбеддингами: SmolLM2-1.7B (24 слоя,
$d=2048$, $I=8192$, 32 головы MHA) и Llama-3.2-1B (16 слоёв, $d=2048$,
$I=8192$, GQA 32/8). На них повторён полный протокол: профиль, таблица
при $k=k_{\mathrm{sink}}$, кривые по $k$, карта, обрыв, ловушка, нейроны
стока, CATS, планировщик и механизм MLP.

\paragraph{Данные.} Все тексты для калибровки и оценки "--- набор
WikiText~\cite{merity2017pointer}: статьи англоязычной Википедии категорий
Good и Featured, собранные и опубликованные Salesforce Research
(вариант \texttt{raw}, без замены редких слов; HF \texttt{Salesforce/wikitext}).
Калибровка "--- WikiText-103 train (103~млн слов): из первых $\approx$2.1~млн символов
корпуса с фиксированным seed берутся 256 случайных неперекрывающихся окон по 512 токенов
(131\,072 токена; для табличных методов 1--2~млн), окна вырезаются из сплошного
текста без начального токена. Оценка "--- WikiText-2 test ($\approx$245~тыс. слов):
первые 32 неперекрывающихся окна по 1024 токена (32\,736 предсказаний).
Обучающая часть WikiText-103 и тест WikiText-2 не пересекаются. Токенизация "---
собственный токенизатор каждой модели (Qwen3: 151\,936 токенов). Для проверки
переноса на другой домен используется pile-10k "--- 10~тыс. документов из The
Pile~\cite{gao2020pile}. Бенчмарки берутся из lm-evaluation-harness~\cite{gao2023harness}
в стандартной постановке (ранжирование вариантов ответа по логарифму правдоподобия):
LAMBADA~\cite{paperno2016lambada} (5\,153 отрывка), HellaSwag~\cite{zellers2019hellaswag}
(10\,042), ARC-Easy/Challenge~\cite{clark2018arc} (2\,376/1\,172), PIQA~\cite{bisk2020piqa}
(1\,838), WinoGrande~\cite{sakaguchi2020winogrande} (1\,267) и MMLU~\cite{hendrycks2021mmlu}
(14\,042 вопроса по 57 предметам); на ноутбуке используются подвыборки (см. ниже).

\paragraph{Метрики.} Основная "--- средняя на токен
$\mathrm{KL}(p_f\,\|\,p_{\tilde f})$ между исходной и изменённой моделью,
критерий <<KL при подстановке>>~\cite{gao_notclean}. Дополнительно:
perplexity, доля совпадения наиболее вероятного токена (top-1) и
относительная ошибка остаточного потока после слоя $k$. Обе модели
считаются в одном процессе в fp32: в bf16 расхождение двух реализаций
одного и того же внимания уже даёт KL$\approx10^{-3}$, а в fp32
вырожденные режимы всех методов дают KL$<10^{-8}$. Бенчмарки
(lm-eval 0.4.13, bf16): LAMBADA, HellaSwag, ARC-e, ARC-c, PIQA,
WinoGrande (0-shot, по 200 примеров) и MMLU (5-shot, по 5 вопросов на
подзадачу, 285 вопросов).

\paragraph{Глубина замены.} $k\in\{1,2,3,4,6\}$ для развёртки по
глубине на 0.6B. Для итоговых таблиц используется $k=k_{\mathrm{sink}}$:
число первых слоёв до появления массивной активации первого токена
(норма больше медианы в 100 раз). Эта глубина определяется
автоматически по профилю каждой модели.

\paragraph{Бюджет.} Для MLP указывается доля параметров и доля
умножений на токен относительно исходного $3dI$. Для внимания "---
доля голов или позиций, для которых вычисляется $q^\top k$. Наши
реализации эталонные (материализуют матрицу $T\times T$), поэтому
время работы не измеряется: сравнение проводится по асимптотике и
числу операций.

\paragraph{Вычисления и воспроизводимость.} Все вычисления выполнены
на ноутбуке с Apple M4 (16 ГБ, MPS). Модели 4B и 8B в fp32 в эту память не
помещаются, поэтому исследование ограничено 0.6B и 1.7B. PyTorch 2.15,
transformers 5.22, lm-eval 0.4.13. Все случайные выборки (калибровочные
окна, инициализация $k$-means, признаки Performer) фиксированы seed. Код,
конфигурации и сырые JSON-результаты прилагаются.
